\chapter{Aspectos éticos, sociales, económicos y ambientales}
\label{cap:aspectos}

\section{Introducción}

Un sistema de detección de intrusiones opera necesariamente sobre registros de la actividad
de personas y sistemas, y produce salidas que orientan decisiones con consecuencias
operativas. Ambas circunstancias exigen una reflexión explícita sobre las implicaciones
éticas y legales del trabajo, que se aborda en este capítulo.

\section{Protección de datos personales}

La telemetría de seguridad contiene de forma inherente información susceptible de constituir
dato personal en el sentido del Reglamento General de Protección de Datos: nombres de
usuario, identificadores de cuenta, nombres de equipo, direcciones IP y, de forma indirecta,
el patrón de actividad de personas identificables.

En el ámbito de este trabajo, la exposición se limita de tres maneras. En primer lugar, los
datos empleados proceden exclusivamente de \textbf{conjuntos públicos capturados en entornos
de laboratorio} construidos al efecto, en los que las cuentas y equipos son ficticios y no
corresponden a personas reales. En segundo lugar, no se ha procesado telemetría de ningún
sistema en producción, propio o de terceros. Y en tercer lugar, los datos permanecen
localmente en el entorno de desarrollo, sin transferencia a servicios externos.

De cara a un eventual despliegue en un entorno real, el sistema plantearía las obligaciones
habituales de esta categoría de tratamiento: existencia de base jurídica —típicamente el
interés legítimo en la seguridad de las redes y la información—, información previa a las
personas afectadas, limitación de la finalidad al ámbito estrictamente de seguridad,
minimización de los datos recogidos, establecimiento de plazos de conservación proporcionados
y control de acceso a la información. El diseño del sistema facilita varias de estas
obligaciones: el particionado temporal del almacén permite implementar políticas de
conservación por rangos, y la separación entre campos promocionados y evento completo
posibilitaría aplicar seudonimización selectiva.

Merece mención específica el principio de \textbf{minimización} frente a la tentación
contraria propia del dominio: la práctica habitual en detección es recoger la máxima
telemetría posible por si resultara necesaria durante una investigación. Existe aquí una
tensión real entre eficacia de detección y proporcionalidad del tratamiento, que en un
despliegue real debería resolverse de forma explícita y documentada, no por omisión.

\section{Uso responsable del conocimiento ofensivo}

El trabajo emplea conocimiento sobre técnicas adversarias reales: cómo se extraen
credenciales de la memoria, cómo se progresa lateralmente por una red o cómo se eliminan las
trazas de una intrusión. Este conocimiento es intrínsecamente de doble uso.

Tres consideraciones acotan el riesgo. La primera es que todo el material empleado procede de
\textbf{fuentes públicas y ampliamente documentadas}: MITRE ATT\&CK, los conjuntos de datos
abiertos utilizados y el repositorio comunitario de reglas de detección son recursos
publicados precisamente con finalidad defensiva. El trabajo no genera conocimiento ofensivo
nuevo ni descubre técnicas no documentadas. La segunda es que la orientación del sistema es
exclusivamente \textbf{detectiva}: el objetivo es reconocer estas técnicas cuando ocurren, no
facilitar su ejecución. La tercera es que la emulación de comportamiento adversario, prevista
en fases posteriores, se realizará únicamente sobre entornos de laboratorio aislados y
controlados.

\section{Sesgo, fiabilidad y supervisión humana}

Un sistema que prioriza incidentes influye directamente en qué recibe atención y qué no. Un
sesgo en la priorización no produce un error visible: produce la \textbf{ausencia silenciosa}
de una alerta que sí debería haberse atendido. Esta consideración ha determinado dos
decisiones de diseño con dimensión ética además de técnica.

La primera es la \textbf{explicabilidad} como requisito, no como añadido. Una puntuación
opaca traslada al analista la disyuntiva de confiar ciegamente o desconfiar por completo;
una puntuación descompuesta en factores interpretables le permite evaluar críticamente el
juicio del sistema y detectar cuándo su razonamiento no aplica al caso concreto.

La segunda es el carácter de \textbf{soporte a la decisión, y no de decisión automatizada}.
El sistema no ejecuta acciones de respuesta: no aísla equipos, no bloquea cuentas ni
interrumpe conexiones. Ordena y justifica, mientras que la decisión permanece en el analista.
Esta delimitación es deliberada y se considera apropiada para el estado actual de madurez de
este tipo de sistemas.

Debe reconocerse asimismo una limitación inherente al enfoque de evaluación adoptado: los
conjuntos de datos empleados provienen de laboratorios y de campañas concretas, por lo que
los resultados obtenidos no son directamente extrapolables a la diversidad de un entorno de
producción real. Esta limitación se explicita en el análisis de resultados y en las
conclusiones.

\section{Licencias y propiedad intelectual}

El trabajo se apoya íntegramente en componentes de código abierto y conjuntos de datos
públicos, cuyas condiciones de uso se han verificado y respetado. Los conjuntos de datos
empleados se distribuyen bajo licencias que permiten su uso con fines de investigación, y las
bases de conocimiento y formatos utilizados —MITRE ATT\&CK, OCSF, STIX, Sigma— son recursos
abiertos publicados para su adopción por la comunidad.

% TODO: completar con la tabla de licencias concreta de cada componente y dataset empleado
% al cierre del desarrollo.

\section{Impacto social y económico}

El impacto potencial más directo del trabajo se sitúa en el ámbito de las organizaciones que
carecen de capacidad para operar un centro de operaciones de seguridad dotado. Las
herramientas comerciales que proporcionan capacidades equivalentes tienen un coste que las
sitúa fuera del alcance de una parte considerable del tejido empresarial. Una arquitectura
construida íntegramente sobre componentes abiertos, desplegable sobre recursos modestos y
que reduce el volumen de alertas a revisar, reduce la barrera de entrada a este tipo de
capacidades.

Debe señalarse, no obstante, que la reducción de la barrera tecnológica no elimina la
necesidad de criterio experto: el sistema ordena y contextualiza, pero la interpretación y la
respuesta siguen requiriendo conocimiento especializado.

\section{Impacto ambiental}

El procesamiento continuo de telemetría tiene un coste energético asociado. Varias decisiones
de diseño adoptadas por motivos de eficiencia operativa tienen también efecto en este plano:
la selección de componentes de bajo consumo de recursos, la organización del despliegue en
perfiles que evitan mantener servicios activos innecesariamente, la aplicación de
codificaciones de compresión específicas por columna en el almacén analítico, y la
restricción explícita del aprendizaje automático a modos de inferencia sobre modelos
pequeños, descartando el entrenamiento de modelos de gran tamaño.

Esta última decisión, adoptada inicialmente por limitación de recursos, resulta relevante
también desde esta perspectiva: el coste energético del entrenamiento de modelos grandes es
considerable, y el trabajo demuestra que el problema abordado no lo requiere.

El desarrollo se ha realizado íntegramente sobre una única estación de trabajo, sin recurso a
infraestructura de centro de datos.
